\documentclass[12pt, a4paper]{article}
\usepackage{../../../myconfig} % Gọi file cấu hình từ ngoài gốc vào

% ==========================================
% BẮT ĐẦU NỘI DUNG TÀI LIỆU
% ==========================================
\begin{document}

% Truyền 3 tham số: {Nhãn cam} {Chữ to chính giữa} {Chữ ở góc phải Header}
\titlebox{Chủ đề: Hàm số}{Hàm logarit, mũ}{Hàm số: Logarit, mũ}

% --- CÂU 1 ---
\questionLinesManual{5}{1}{Rút gọn biểu thức \( P = x^{\frac{1}{3}} \sqrt{x} \) với \( x > 0 \).}{
    \begin{tasks}(4)
        \task \( P = \sqrt{x} \)
        \task \( P = x^{\frac{1}{8}} \)
        \task \( P = x^{\frac{2}{9}} \)
        \task \( P = x^2 \)
    \end{tasks}
}

% --- CÂU 2 ---
\questionLinesManual{5}{2}{Với \( a > 0 \), \( b > 0 \), \( \alpha, \beta \) là các số thực bất kì, đẳng thức nào sau đây sai?}{
    \begin{tasks}(4)
        \task \( \dfrac{a^{\alpha}}{a^{\beta}} = a^{\alpha - \beta} \)
        \task \( a^{\alpha} \cdot a^{\beta} = a^{\alpha + \beta} \)
        \task \( \dfrac{a^{\alpha}}{b^{\beta}} = \left( \dfrac{a}{b} \right)^{\alpha - \beta} \)
        \task \( a^{\alpha} \cdot b^{\alpha} = (ab)^{\alpha} \)
    \end{tasks}
}

% --- CÂU 3 ---
\questionLinesManual{5}{3}{Cho \( a > 0; a \neq 1 \). Tìm mệnh đề đúng trong các mệnh đề sau?}{
    \begin{tasks}(2)
        \task \( \log_a a = 1 \)
        \task \( \log_a x \) có nghĩa \( \forall x \in \mathbb{R} \)
        \task \( \log_a a = 0 \)
        \task \( \log_a (x \cdot y) = \log_a x \cdot \log_a y; \forall x, y > 0 \)
    \end{tasks}
}

% --- CÂU 4 ---
\questionLinesManual{5}{4}{Với \( a, b \) là các số thực dương tùy ý và \( a \neq 1 \), \( \log_a \left( \sqrt[3]{b} \right) \) bằng}{
    \begin{tasks}(4)
        \task \( 3 + \log_a b \)
        \task \( 3 \log_a b \)
        \task \( \dfrac{1}{3} + \log_a b \)
        \task \( \dfrac{1}{3} \log_a b \)
    \end{tasks}
}

% --- CÂU 5 ---
\questionLinesManual{5}{5}{Cho \( a \) và \( b \) là hai số thực dương thoả mãn \( a^3 b^2 = 32 \). Giá trị của \( 3 \log_2 a + 2 \log_2 b \) bằng}{
    \begin{tasks}(4)
        \task $5$
        \task $2$
        \task $32$
        \task $4$
    \end{tasks}
}

% --- CÂU 6 ---
\questionLinesManual{5}{6}{Hàm số nào sau đây đồng biến trên \( \mathbb{R} \)?}{
    \begin{tasks}(4)
        \task \( y = \ln x \)
        \task \( y = 3^x \)
        \task \( y = \left( \dfrac{1}{2} \right)^x \)
        \task \( y = \log_2 x \)
    \end{tasks}
}

% --- CÂU 7 ---
\questionLinesManual{5}{7}{Tập xác định của hàm số \( y = \log_6 x \) là}{
    \begin{tasks}(4)
        \task \([0; +\infty)\)
        \task \((0; +\infty)\)
        \task \((-\infty; 0)\)
        \task \((-\infty; +\infty)\)
    \end{tasks}
}

% --- CÂU 8 ---
\questionLinesManual{5}{8}{Nghiệm của phương trình \( 2^{2x-1} = 32 \) là}{
    \begin{tasks}(4)
        \task $x = 3$
        \task $x = \dfrac{17}{2}$
        \task $x = \dfrac{5}{2}$
        \task $x = 2$
    \end{tasks}
}

% --- CÂU 9 ---
\questionLinesManual{5}{9}{Nghiệm của phương trình \( \log_5(3x-5) = \log_5(5-x) \) là:}{
    \begin{tasks}(4)
        \task $x = 2$
        \task $x = \dfrac{2}{5}$
        \task $x = 3$
        \task $x = \dfrac{5}{2}$
    \end{tasks}
}

% --- CÂU 10 ---
\questionLinesManual{5}{10}{Tập nghiệm của bất phương trình \( 10^{2x+3} \leq \dfrac{1}{10} \) là}{
    \begin{tasks}(4)
        \task $(-2; +\infty)$
        \task $(-\infty; -2)$
        \task $(-\infty; -2]$
        \task $[-2; +\infty)$
    \end{tasks}
}

% --- CÂU 11 ---
\questionLinesManual{5}{11}{Tập nghiệm của bất phương trình \( \log_7(x-3) \leq 1 \) là:}{
    \begin{tasks}(4)
        \task $[3; 10]$
        \task $(3; 10]$
        \task $(3; 10)$
        \task $[3; 10)$
    \end{tasks}
}

% --- CÂU 12 ---
\questionLinesManual{5}{12}{Số nghiệm nguyên của bất phương trình \( \log_3(x+2) - \log_3(x-2) \leq 3 \) là}{
    \begin{tasks}(4)
        \task $3$
        \task $2$
        \task $4$
        \task $5$
    \end{tasks}
}

% --- CÂU 13 ---
\questionLinesManual{5}{13}{Tính giá trị của biểu thức \( E = 9^{2-\sqrt{3}} \cdot 3^{-1+2\sqrt{3}} \).}{
    \begin{tasks}(4)
        \task $E = 9$
        \task $E = \dfrac{1}{3}$
        \task $E = 27$
        \task $E = 3$
    \end{tasks}
}

% --- CÂU 14 ---
\questionLinesManual{5}{14}{Tính \( \log_2 16 \).}{
    \begin{tasks}(4)
        \task $2$
        \task $4$
        \task $8$
        \task $16$
    \end{tasks}
}

% --- CÂU 15 ---
\questionLinesManual{5}{15}{Giá trị của \( \log_2 \dfrac{1}{16} \) bằng}{
    \begin{tasks}(4)
        \task $4$
        \task $\dfrac{1}{4}$
        \task $\dfrac{1}{8}$
        \task $-4$
    \end{tasks}
}

% --- CÂU 16 ---
\questionLinesManual{5}{16}{Tập xác định của hàm số \( y = \log_3 2x \) là}{
    \begin{tasks}(4)
        \task $(-\infty; 0)$
        \task $(0; +\infty)$
        \task $\mathbb{R}$
        \task $(1; +\infty)$
    \end{tasks}
}

% --- CÂU 17 ---
\questionLinesManual{5}{17}{Nghiệm của phương trình \( 5^x = 2 \) là:}{
    \begin{tasks}(4)
        \task \( x = \log_2 5 \)
        \task \( x = \log_5 2 \)
        \task \( x = \dfrac{2}{5} \)
        \task \( x = \sqrt{5} \)
    \end{tasks}
}

% --- CÂU 18 ---
\questionLinesManual{5}{18}{Tìm tập nghiệm của bất phương trình \( 0,4^{x^2+x} > 0,16 \).}{
    \begin{tasks}(4)
        \task $(-2;1)$
        \task $[-2;1]$
        \task $(-\infty;-2)$
        \task $(1;+\infty)$
    \end{tasks}
}

% --- CÂU 19 ---
\questionLinesManual{5}{19}{Tập nghiệm của phương trình \( \log_3 (3x-1)^2 = 2 \) là}{
    \begin{tasks}(4)
        \task \( \left\{ \dfrac{4}{3} \right\} \)
        \task \( \left\{ \dfrac{4}{3}; -\dfrac{2}{3} \right\} \)
        \task \( \left\{ -\dfrac{2}{3} \right\} \)
        \task \( \left\{ \dfrac{10}{3}; -\dfrac{2}{3} \right\} \)
    \end{tasks}
}

% --- CÂU 20 ---
\questionLinesManual{5}{20}{Số nghiệm nguyên của bất phương trình \( \log_3^2 (2x-1) - 3 \cdot \log_3 (2x-1)^2 + 2 < 0 \) là}{
    \begin{tasks}(4)
        \task $1$
        \task $2$
        \task $0$
        \task $3$
    \end{tasks}
}

% --- CÂU 21 ---
\questionLinesManual{5}{21}{Tập nghiệm của bất phương trình: \( 3^{\sqrt{x+1}-2} \leq 0 \) là}{
    \begin{tasks}(4)
        \task $[-1;3)$
        \task $[-1;+\infty)$
        \task $[-1;3]$
        \task $(\infty;3)$
    \end{tasks}
}

% --- CÂU 22 ---
\questionLinesManual{5}{22}{Biểu thức \( T = \sqrt[3]{a^3} \) với \( a > 0 \), được viết dưới dạng luỹ thừa với số mũ hữu tỉ là}{
    \begin{tasks}(4)
        \task \( a^{\frac{3}{5}} \)
        \task \( a^{\frac{2}{15}} \)
        \task \( a^{\frac{4}{15}} \)
        \task \( a^{\frac{3}{10}} \)
    \end{tasks}
}

% --- CÂU 23 ---
\questionLinesManual{5}{23}{Rút gọn biểu thức \( P = x^{\frac{1}{2}} \cdot \sqrt{x} \) (với \( x > 0 \))}{
    \begin{tasks}(4)
        \task \( x^4 \)
        \task \( x^{\frac{1}{16}} \)
        \task \( x^{\frac{5}{16}} \)
        \task \( x^8 \)
    \end{tasks}
}

% --- CÂU 24 ---
\questionLinesManual{5}{24}{Cho \( \log_2 a = 5, \log_2 b = 7; a, b > 0 \). Giá trị của \( \log_2 (a^2 b^3) \) bằng}{
    \begin{tasks}(4)
        \task $12$
        \task $31$
        \task $35$
        \task $13$
    \end{tasks}
}

% --- CÂU 25 ---
\questionLinesManual{5}{25}{Cho \( 0 < a \neq 1, \log_3 5 = a \). Giá trị của \( \log_5 75 \) theo \( a \) bằng}{
    \begin{tasks}(4)
        \task $2a$
        \task $2 + \dfrac{1}{a}$
        \task $2 - \dfrac{1}{a}$
        \task $2 + a$
    \end{tasks}
}

% --- CÂU 26 ---
\questionLinesManual{5}{26}{Trong các hàm số sau đây, hàm nào là hàm số mũ?}{
    \begin{tasks}(4)
        \task \( y = \sqrt{2^x} \)
        \task \( y = x^3 \)
        \task \( y = -3x^2 \)
        \task \( y = 5^{\frac{1}{3}} \)
    \end{tasks}
}

% --- CÂU 27 ---
\questionLinesManual{5}{27}{Tập xác định của hàm số \( y = \log_7 (x - 1) \) là}{
    \begin{tasks}(4)
        \task \( \mathbb{R} \setminus \{1\} \)
        \task \( [1; +\infty) \)
        \task \( \mathbb{R} \)
        \task \( (1; +\infty) \)
    \end{tasks}
}

% --- CÂU 28 ---
\questionLinesManual{5}{28}{Nghiệm của phương trình \( \log_2 (2x) = 3 \) là:}{
    \begin{tasks}(4)
        \task $x = 4$
        \task $x = 3$
        \task $x = 1$
        \task $x = 2$
    \end{tasks}
}

% --- CÂU 29 ---
\questionLinesManual{5}{29}{Tập nghiệm của bất phương trình \( 3^x \leq 184 \) là}{
    \begin{tasks}(4)
        \task \( S = (-\infty; \log_3 184] \)
        \task \( S = [\log_3 184; +\infty) \)
        \task \( S = (-\infty; \log_3 184) \)
        \task \( S = (\log_3 184; +\infty) \)
    \end{tasks}
}

% --- CÂU 30 ---
\questionLinesManual{5}{30}{Tìm nghiệm của phương trình: \( 4^{-5x-10} = 16^{6-10x} \)}{
    \begin{tasks}(4)
        \task \( x = \dfrac{22}{5} \)
        \task \( x = \dfrac{97}{15} \)
        \task \( x = \dfrac{16}{5} \)
        \task \( x = \dfrac{22}{15} \)
    \end{tasks}
}

% --- CÂU 31 ---
\questionLinesManual{5}{31}{Tập nghiệm của bất phương trình \( \log_4 x > -1 \).}{
    \begin{tasks}(4)
        \task \( S = \left( -\infty; \dfrac{1}{4} \right) \)
        \task \( S = (-\infty; 1) \)
        \task \( S = \left( \dfrac{1}{4}; +\infty \right) \)
        \task \( S = (1; +\infty) \)
    \end{tasks}
}

% --- CÂU 32 ---
\questionLinesManual{5}{32}{Với \( a \) là số thực dương tùy ý, biểu thức \( a^3 \cdot a^3 \) là}{
    \begin{tasks}(4)
        \task \( a^2 \)
        \task \( a^5 \)
        \task \( a^{\frac{5}{9}} \)
        \task \( a^{\frac{4}{3}} \)
    \end{tasks}
}

% --- CÂU 33 ---
\questionLinesManual{5}{33}{Khẳng định nào sau đây đúng?}{
    \begin{tasks}(4)
        \task \( 3^{\frac{1}{3}} = \dfrac{1}{3^3} \)
        \task \( 3^{\frac{1}{16}} = \dfrac{3}{3^{16}} \)
        \task \( 3^{-1} = -3 \)
        \task \( 3^{-0,3} = \dfrac{1}{3^{0,3}} \)
    \end{tasks}
}

% --- CÂU 34 ---
\questionLinesManual{5}{34}{Cho \( x, y > 0 \) và \( \alpha, \beta \in \mathbb{R} \). Tìm đẳng thức sai dưới đây.}{
    \begin{tasks}(4)
        \task \( x^\alpha + y^\alpha = (x + y)^\alpha \)
        \task \( x^\alpha \cdot y^\alpha = (xy)^\alpha \)
        \task \( x^\alpha \cdot x^\beta = x^{\alpha + \beta} \)
        \task \( (x^\alpha)^\beta = x^{\alpha \beta} \)
    \end{tasks}
}

% --- CÂU 35 ---
\questionLinesManual{5}{35}{Với \( 0 < a \neq 1, M > 0; \alpha \in \mathbb{R} \). Tìm mệnh đề sai trong các mệnh đề sau}{
    \begin{tasks}(4)
        \task \( \log_a 1 = 0 \)
        \task \( \log_a 1 = 1 \)
        \task \( \log_a a^\alpha = \alpha \)
        \task \( a^{\log_a M} = M \)
    \end{tasks}
}

% --- CÂU 36 ---
\questionLinesManual{5}{36}{Cho \( \log_2 a = 5, a > 0 \). Giá trị của \( \log_2 2a \) bằng}{
    \begin{tasks}(4)
        \task $7$
        \task $6$
        \task $-4$
        \task $10$
    \end{tasks}
}

% --- CÂU 37 ---
\questionLinesManual{5}{37}{Cho \( 0 < a \neq 1, \log_3 a = 2 \). Giá trị của \( \log_a 3a \) bằng}{
    \begin{tasks}(4)
        \task $7$
        \task \( \dfrac{3}{2} \)
        \task $6$
        \task $3$
    \end{tasks}
}

% --- CÂU 38 ---
\questionLinesManual{5}{38}{Trong các hàm số sau đây, hàm số nào nghịch biến trên \( \mathbb{R} \)?}{
    \begin{tasks}(4)
        \task \( y = \left( \dfrac{1}{3} \right)^{-x} \)
        \task \( y = \left( \sqrt{2} \right)^x \)
        \task \( y = \left( \sqrt{3} \right)^x \)
        \task \( y = \left( \dfrac{1}{2} \right)^x \)
    \end{tasks}
}

% --- CÂU 39 ---
\questionLinesManual{5}{39}{Nghiệm của phương trình \( \log_2(x+1) = 3 \) là}{
    \begin{tasks}(4)
        \task $x = 2$
        \task $x = 8$
        \task $x = 1$
        \task $x = 7$
    \end{tasks}
}

% --- CÂU 40 ---
\questionLinesManual{5}{40}{Tập nghiệm của bất phương trình \( 5^x < 70 \) là}{
    \begin{tasks}(4)
        \task \( S = [\log_5 70; +\infty) \)
        \task \( S = (\log_5 70; +\infty) \)
        \task \( S = (-\infty; \log_5 70) \)
        \task \( S = (-\infty; \log_5 70] \)
    \end{tasks}
}

% --- CÂU 41 ---
\questionLinesManual{5}{41}{Tìm nghiệm của phương trình \( 4^{x-1} = \dfrac{1}{16} \).}{
    \begin{tasks}(4)
        \task $x = 5$
        \task $x = 4$
        \task $x = -1$
        \task $x = -10$
    \end{tasks}
}

% --- CÂU 42 ---
\questionLinesManual{5}{42}{Tập nghiệm của bất phương trình \( \log_4 x < -2 \).}{
    \begin{tasks}(4)
        \task \( S = (-\infty; 16) \)
        \task \( S = \left( 0; \dfrac{1}{16} \right) \)
        \task \( S = \left( \dfrac{1}{16}; +\infty \right) \)
        \task \( S = (16; +\infty) \)
    \end{tasks}
}

% --- CÂU 43 ---
\questionLinesManual{5}{43}{Tập nghiệm của bất phương trình \( \log x \geq 1 \) là}{
    \begin{tasks}(4)
        \task \( (10; +\infty) \)
        \task \( (0; +\infty) \)
        \task \( [10; +\infty) \)
        \task \( (-\infty; 10) \)
    \end{tasks}
}

% --- CÂU 44 ---
\questionLinesManual{5}{44}{Tập nghiệm của bất phương trình \( \left( \dfrac{1}{3} \right)^x > 9 \) là}{
    \begin{tasks}(4)
        \task \( (2; +\infty) \)
        \task \( (-\infty; -2) \)
        \task \( (-\infty; 2) \)
        \task \( (-2; +\infty) \)
    \end{tasks}
}

% --- CÂU 45 ---
\questionLinesManual{5}{45}{Tập nghiệm của bất phương trình \( 2^x > 1 \) là}{
    \begin{tasks}(4)
        \task \( (0; +\infty) \)
        \task \( (-\infty; 0) \)
        \task \( (0; 1) \)
        \task \( (0; 4) \)
    \end{tasks}
}

% --- CÂU 46 ---
\questionLinesManual{5}{46}{Tập nghiệm của bất phương trình \( \log_3 (x - 2) > 2 \) là}{
    \begin{tasks}(4)
        \task \( (2; +\infty) \)
        \task \( (2; 11) \)
        \task \( (1; +\infty) \)
        \task \( (-1; 2) \)
    \end{tasks}
}

% --- CÂU 47 ---
\questionLinesManual{5}{47}{Tập nghiệm của bất phương trình \( \log_{0,2}(x-1) < 0 \) là}{
    \begin{tasks}(4)
        \task \( (1; 2) \)
        \task \( (2; +\infty) \)
        \task \( (-\infty; 1) \)
        \task \( (-\infty; 2) \)
    \end{tasks}
}

% --- CÂU 48 ---
\questionLinesManual{5}{48}{Tập nghiệm của bất phương trình \( \log_5(2x-1) < \log_5(x+2) \) là}{
    \begin{tasks}(4)
        \task \( S = (3; +\infty) \)
        \task \( S = (-2; 3) \)
        \task \( S = \left( \dfrac{1}{2}; 3 \right) \)
        \task \( S = (-\infty; 3) \)
    \end{tasks}
}

% --- CÂU 49 ---
\questionLinesManual{5}{49}{Số nghiệm nguyên của bất phương trình \( \log_{0,5}(2x+6) \geq -5 \) là}{
    \begin{tasks}(4)
        \task $16$
        \task $13$
        \task $15$
        \task $8$
    \end{tasks}
}

% --- CÂU 50 ---
\questionLinesManual{5}{50}{Tìm tập xác định \( D \) của hàm số \( y = \log_2 (x^2 - 2x - 3) \).}{
    \begin{tasks}(2)
        \task \( D = (-\infty; -1) \cup (3; +\infty) \)
        \task \( D = (-\infty; -1] \cup [3; +\infty) \)
        \task \( D = [-1; 3] \)
        \task \( D = (-1; 3) \)
    \end{tasks}
}

\end{document}